\addcontentsline{toc}{subsubsection}{Solution 14. IMO 2023 Shortlist -- Algebra A4}

\begin{dlsummary}[title={\bfseries Solution 14. IMO 2023 Shortlist -- Algebra A4},breakable]

\begin{proof}
Write $f^n(x)$ for the $n$-fold composition
of $f$, with $f^0(x)=x$.

\medskip
\noindent\textbf{Step 1.}
Setting $y=x$ gives
\[
2xf(x)\ge(f^2(x)+x)f(x).
\]
Since $f(x)>0$,
\[
\boxed{f^2(x)\le x.} \tag{1}
\]

\medskip
\noindent\textbf{Step 2.}
Setting $x=f(y)$ yields
\[
f(y)-f^3(y)\ge y-f^2(y).
\]
Replacing $y$ by $f^n(y)$ gives
\[
f^{n+1}(y)-f^{n+3}(y)
\ge f^n(y)-f^{n+2}(y).
\]

Define
\[
D_n=f^n(y)-f^{n+2}(y).
\]
Then
\[
D_{n+1}\ge D_n\ge D_0
=y-f^2(y)\ge0.
\]

For every positive integer $m$,
\begin{align*}
y-f^{2m}(y)
&=\sum_{k=0}^{m-1}D_{2k}\\
&\ge m\bigl(y-f^2(y)\bigr).
\end{align*}

Since $f^{2m}(y)>0$,
\[
y>m\bigl(y-f^2(y)\bigr).
\]
As $m$ is arbitrarily large,
\[
\boxed{f^2(y)=y.} \tag{2}
\]

\medskip
\noindent\textbf{Step 3.}
Using (2), the original inequality becomes
\[
x(f(x)+f(y))\ge(x+y)f(y),
\]
hence
\[
xf(x)\ge yf(y).
\]
Interchanging $x$ and $y$ gives equality:
\[
xf(x)=yf(y)=c
\]
for some $c>0$. Therefore,
\[
\boxed{f(x)=\frac{c}{x},\qquad c>0.}
\]

Conversely, for $f(x)=c/x$,
we have $f(f(x))=x$ and
\[
x(f(x)+f(y))
=c+\frac{cx}{y}
=(f(f(x))+y)f(y).
\]
Thus every such function satisfies the
given inequality.
\end{proof}

\end{dlsummary}